% r03_hyb_ES: proof of \Cref{rot:prop:BD}, part VI: the continuum extension of the Y kernels (cubic-Hermite cardinal interpolation), its decay, its multiplier form at the X sites, and the decaying forms at the X sites.

\paragraph{The continuum extension of the $Y$ kernels.}
The comparisons of \Cref{rot:hyb:lem:CV,rot:hyb:lem:CLN} below compare a sum over $X$ bonds with a sum over $Y$ bonds through \Cref{rot:hyb:lem:CL}, which needs a $C^1$ field on the block that equals the $Y$ kernel on the $Y$ lattice. The Fourier formula of \Cref{rot:lem:H} evaluated off the lattice (the sharp-zone extension) is not admissible for this: the symbol does not vanish on the faces of the zone, so the extension and its gradient have algebraic tails of order $\eta'^3/|r-z|$ and $\eta'^2/|r-z|$ off the lattice, not summable over a torus in $d=3$. The field used is the cubic-Hermite cardinal interpolation of the lattice kernels: it is $C^1$, equals the lattice kernel on $\Lambda_Y$, has exponentially decaying value and gradient by the lattice bounds alone, and obeys the multiplier comparison of \Cref{rot:lem:H} at the $X$ sites with an explicit constant.

\begin{definition}[cubic-Hermite cardinal extension]\label{rot:hyb:def:ES}
$S(u)\defeq3u^2-2u^3$, $\psi(s)\defeq S(1-|s|)\mathbf 1_{|s|<1}=(1-|s|)^2(1+2|s|)\mathbf 1_{|s|<1}$, $\Psi(v)\defeq\prod_\alpha\psi(v\cdot f_\alpha/\eta')$ (the frame of $\Yk$). For $F$ on $\Lambda_Y=R^{-1}(\eta'\Z^3)$ put $(I_SF)(r)\defeq\sum_{y\in\Lambda_Y}\Psi(r-y)F(y)$, $r\in\R^3$. The extensions are $E^S(r,z)\defeq I_S[H^Y(\cdot,z)](r)$ for the level-$r$ lift (scalar and photon), $E^S_j(r,z)\defeq I_S[H^Y_j(\cdot,z)](r)$ for the level-$j$ lift, $\tilde S^S_j(r,y)\defeq\sum_{z,z'}E^S_j(r,z)C^Y_j(z,z')\overline{H^Y_j(y,z')}=I_S[S^Y_j(\cdot,y)](r)$ for $y\in\Lambda_Y$ ($S^Y_j=H^Y_jC^Y_jH^{Y*}_j$, \Cref{rot:def:match}), likewise for the photon slices, and $E^S\varphi\defeq I_S[H^Y\varphi]$ for a coarse field $\varphi$ (componentwise in any frame: $I_S$ is scalar and linear). The extension acts in the first argument only; second arguments stay on $\Lambda_Y$.
\end{definition}

\begin{lemma}[cell identity]\label{rot:hyb:lem:ESH}
Write the frame coordinates of $r$ as $f_\alpha\cdot r=\eta'(n_\alpha+u_\alpha)$, $n\in\Z^3$, $u\in[0,1)^3$, with cell corners $y_\sigma\defeq R^{-1}\eta'(n+\sigma)$, $\sigma\in\{0,1\}^3$, and weights $\omega_\sigma(u)\defeq\prod_\alpha\omega_{\sigma_\alpha}(u_\alpha)$, $\omega_0\defeq1-S$, $\omega_1\defeq S$ (so $\omega_\sigma\ge0$, $\sum_\sigma\omega_\sigma=1$). Then
\textup{(i)} $I_SF(r)=\sum_\sigma\omega_\sigma(u)F(y_\sigma)$;
\textup{(ii)} $\partial_{f_1}I_SF(r)=S'(u_1)\sum_{\sigma'\in\{0,1\}^2}\omega_{\sigma'}(u_2,u_3)\,(\partial_1F)(y_{(0,\sigma')})$ with the lattice forward difference $\partial_1F(y)\defeq[F(y+\eta'f_1)-F(y)]/\eta'$, and likewise for $f_2,f_3$;
\textup{(iii)} $|I_SF(r)|\le\max_\sigma|F(y_\sigma)|$ and $|\partial_{f_\alpha}I_SF(r)|\le\frac32\max|\partial_\alpha F|$ over the four cell edges parallel to $f_\alpha$, so $|\nabla I_SF(r)|\le\frac{3\sqrt3}2\max_{\mathrm{edges}}|\partial F|$;
\textup{(iv)} $I_SF\in C^1(\R^3)$ and $I_SF=F$ on $\Lambda_Y$.
\end{lemma}
\begin{proof}
Only the corners with $|(r-y)\cdot f_\alpha|<\eta'$ for all $\alpha$ contribute, i.e.\ the $y_\sigma$; for them $\psi((r-y_\sigma)\cdot f_\alpha/\eta')=\psi(u_\alpha-\sigma_\alpha)$, which is $S(1-u_\alpha)=1-S(u_\alpha)$ for $\sigma_\alpha=0$ and $S(u_\alpha)$ for $\sigma_\alpha=1$ ($S(u)+S(1-u)=1$): this is (i). Differentiating (i) in $r\cdot f_1=\eta'(n_1+u_1)$ gives $\eta'^{-1}S'(u_1)\sum_{\sigma'}\omega_{\sigma'}[F(y_{(1,\sigma')})-F(y_{(0,\sigma')})]$, which is (ii). (iii) follows from $\omega\ge0$, $\sum\omega=1$, $0\le S'(u)=6u(1-u)\le\frac32$. (iv) $\psi'(s)=-6s(1-|s|)$ on $|s|<1$ vanishes at $s=0,\pm1$, so $\psi\in C^1(\R)$, and the sum is locally finite; $\psi(m)=\delta_{m0}$ for $m\in\Z$.
\end{proof}

\begin{lemma}[decay of the extensions]\label{rot:hyb:lem:ESB}
The lattice inputs are \cite[(3.6), (3.7), (3.63)]{K1} for the standard flow: $|H^Y(y,z)|,|\partial H^Y(y,z)|\le Ce^{-\delta_0|y-z|}$ ((3.7) with $f=Q^*\delta_z$, supported in the block of $z$, $\|f\|_\infty\le1$; for the photon, ``the same bounds for the vector field operators'', \cite[p.~656]{K1}); at level $j$, in the unit scale of level $j$, $|H^Y_j|,|\partial H^Y_j|\le Ce^{-\delta_0|\cdot|}$ and $|C^Y_j(z,z')|\le Ce^{-\delta_0|z-z'|}$. Then for $0<\delta<\delta_0$, uniformly in $j,r,\eta'$, the torus and $m_j\ge0$, for all $r\in\R^3$, $z$, and $y\in\Lambda_Y$, in the unit scale of level $r$:
\begin{enumerate}[label=\textup{(\alph*)},leftmargin=*]
\item $|E^S(r,z)|,\ \frac{2}{3\sqrt3}|\nabla_rE^S(r,z)|\le Ce^{\sqrt3\delta_0\eta'}e^{-\delta_0|r-z|}$, and the same for $E^S_j$ in the unit scale of level $j$;
\item $|\tilde S^S_j(r,y)|\le C\ell_j^{-1}e^{-\delta|r-y|/\ell_j}$, $|\nabla_r\tilde S^S_j(r,y)|,\ |\partial_y\tilde S^S_j(r,y)|\le C\ell_j^{-2}e^{-\delta|r-y|/\ell_j}$ and $|\nabla_r\partial_y\tilde S^S_j(r,y)|\le C\ell_j^{-3}e^{-\delta|r-y|/\ell_j}$ ($\partial_y$ a lattice difference in the second argument);
\item $E^S(\cdot,z)$, $\tilde S^S_j(\cdot,y)$, $E^S\varphi$ are $C^1$ and equal $H^Y(\cdot,z)$, $S^Y_j(\cdot,y)$, $H^Y\varphi$ on $\Lambda_Y$.
\end{enumerate}
\end{lemma}
\begin{proof}
(a) \Cref{rot:hyb:lem:ESH}(iii) with $F=H^Y(\cdot,z)$ and $|y_\sigma-r|\le\sqrt3\eta'$. (b) In the unit scale of level $j$ put $F(y')\defeq S^Y_j(y',y)$, resp.\ $\partial_yS^Y_j(y',y)$ ($\partial_y$ acts on the second argument and commutes with $I_S$). With $S^Y_j=H^Y_jC^Y_jH^{Y*}_j$, $F$ and its lattice gradient in $y'$ are sums $\sum_{z,z'}$ of products of three kernels with at most one lattice difference each, so with $e^{-\delta_0(|y'-z|+|z-z'|+|z'-y|)}\le e^{-\delta|y'-y|}e^{-(\delta_0-\delta)(|y'-z|+|z-z'|)}$ they are $\le C^3K_\delta e^{-\delta|y'-y|}$, $K_\delta\defeq(\sum_{v\in\Z^3}e^{-(\delta_0-\delta)(|v|-\sqrt3)})^2$. \Cref{rot:hyb:lem:ESH}(iii) at the spacing $\eta'/\ell_j\le\frac15$ adds $e^{\sqrt3\delta/5}$ and the factor $\frac{3\sqrt3}2$ for the gradient. Rescale by \cite[(2.20)]{K1}: kernel $\ell_j^{-1}$, each derivative $\ell_j^{-1}$. (c) \Cref{rot:hyb:lem:ESH}(iv). Only lattice inputs with at most one derivative occur; no same-argument second derivative of a block-spin kernel is used.
\end{proof}

\begin{lemma}[multiplier form at the $X$ sites]\label{rot:hyb:lem:ADMS}
Unit scale of level $r$; $x\in\Gamma_1$, $s\in\R^3$, $|s|\le\sqrt{13}\eta'$. In the notation of \Cref{rot:lem:H},
\[
E^S(x+s,z)=(2\pi)^{-3}\!\int_{V_\eta}\!dq\,e^{\mathrm iq\cdot(x-z)}u_{r-1}(q)M^Y_S(q;s)A^Y(q),\qquad M^Y_S(q;s)\defeq\sum_{l\in L^Y_2}\frac{w_l(q;s)\,u_Y(q+l)}{\Delta'_Y(q+l)},\qquad
w_l(q;s)\defeq\sum_{L\in\Lambda_Y^*}\chi(q+l+L)e^{\mathrm i(q+l+L)\cdot s},
\]
$\chi(Q)\defeq\prod_\alpha\hat\psi(\eta'Q\cdot f_\alpha)$, $\hat\psi(t)=\int\psi(\sigma)e^{-\mathrm it\sigma}d\sigma=12(2-2\cos t-t\sin t)/t^4$, $\Lambda_Y^*=\frac{2\pi}{\eta'}R^{-1}\Z^3$; and on $V_\eta$: $|w_l|\le W^3\le1.20$ for $l\ne0$, $|w_0(q;s)-e^{\mathrm iq\cdot s}|\le0.106\,\eta|q|$, $|M^Y_S(q;s)-\Delta_R^{-1}|\le207\,\eta|q|^{-1}$. Consequently, for $x'=x+s\in\Lambda_X$ in the tile of $x$, $|H^X(x',z)-E^S(x',z)|\le C^S_EJ_1\eta$ with $C^S_E\defeq380a+0.2a(C_*/20+8.3a)$ ($C_*$ of \Cref{rot:lem:H}), and the same at level $j$ in the unit scale of level $j$.
\end{lemma}
\begin{proof}
Poisson summation over $\Lambda_Y$ ($\hat\Psi(Q)=\eta'^3\chi(Q)$): $\sum_y\Psi(r-y)e^{\mathrm iQ\cdot y}=\sum_L\chi(Q+L)e^{\mathrm i(Q+L)\cdot r}$. Inserted in $H^Y(y,z)=(2\pi)^{-3}\int_{V_{\eta'}}e^{\mathrm iQ\cdot(y-z)}h(Q)dQ$ ($h$ periodic under $\Lambda_Y^*$, $e^{\mathrm iL\cdot z}=1$) this gives $E^S(r,z)=(2\pi)^{-3}\int_{\R^3}e^{\mathrm iQ\cdot(r-z)}\chi(Q)h(Q)dQ$; splitting $\R^3$ into the translates $V_\eta+l+L$ and using $e^{\mathrm i(l+L)\cdot x}=1$ gives the formula. Bounds: $\psi\ge0$ gives $|\hat\psi|\le\hat\psi(0)=1$; $\hat\psi(2\pi n)=0$ for $n\ne0$; for $|\tau|\le\pi/5$ and $n\ne0$, $|2-2\cos\tau-(2\pi n+\tau)\sin\tau|\le\tau^4/12+2\pi|n||\tau|\le1.004\cdot2\pi|n||\tau|$ and $|2\pi n+\tau|\ge0.9\cdot2\pi|n|$, so $|\hat\psi(2\pi n+\tau)|\le0.074|\tau|/|n|^3$ and $T(\tau)\defeq\sum_{n\ne0}|\hat\psi(\tau+2\pi n)|\le2\zeta(3)\cdot0.074|\tau|\le0.18|\tau|$ (below the cruder $0.213|\tau|$ is used); $1-\hat\psi(t)\le t^2/15$ on $[0,\pi]$ (the series $1-t^2/15+t^4/560-\dots$); $W\defeq\sup_t\sum_n|\hat\psi(t+2\pi n)|\le1.06$ (a one-dimensional bound: the sum is $1$ at the integers and its oscillation on a period is controlled by the bound on $T$). For $l\ne0$: $|w_l|\le\prod_\alpha\sum_n|\hat\psi(t_\alpha+2\pi n)|\le W^3$, $t_\alpha\defeq\eta'(q+l)\cdot f_\alpha$. For $l=0$ ($|t_\alpha|\le\pi/5$ on $V_\eta$, $t\defeq\eta'q$): $|w_0-e^{\mathrm iq\cdot s}|\le(1-\chi(q))+[\prod_\alpha(1+T(t_\alpha))-1]\le|t|^2/15+(e^{0.213\sum_\alpha|t_\alpha|}-1)\le(0.073+0.454)|t|\le0.106\,\eta|q|$ (convexity on $\sum|t_\alpha|\le3\pi/5$). Main term: $|w_0u_Y-1|\le0.106\eta|q|+|q||s|+(1-u_Y)\le(0.106+0.722+0.218)\eta|q|$, over $\Delta'_Y\ge\frac4{\pi^2}|q|^2$: $\le2.59\eta|q|^{-1}$; $|\Delta'^{-1}_Y-\Delta_R^{-1}|\le1.12\eta|q|^{-1}$ (proof of \Cref{rot:lem:H}); the $124$ aliases: $\le124\,W^3\eta^2/4\le37.2\eta^2\le203\eta|q|^{-1}$ ($|q|\le\sqrt3\pi/\eta$). So $|M^Y_S-\Delta_R^{-1}|\le207\eta|q|^{-1}$, and with $|M^X-\Delta_R^{-1}|\le172.1\eta|q|^{-1}$ at the same offset $s$ (proof of \Cref{rot:lem:H}), $|M^X-M^Y_S|\le380\eta|q|^{-1}$. Then the multiplier step of \Cref{rot:lem:H} verbatim: $|M^XA^X-M^Y_SA^Y|\le380a\eta|q|^{-1}+0.2a\eta^2A^Y|M^Y_S|\le C^S_E\eta|q|^{-1}$; integrate against $|u_{r-1}|$.
\end{proof}

\begin{corollary}[decaying forms at the $X$ sites]\label{rot:hyb:cor:DF}
Unit scale of level $r$; $x'\in\Lambda_X$ in the tile of $x\in\Gamma_1$, $x''=\tau x'$; $y'\in\Lambda_X$, $y''=\tau y'$.
\begin{enumerate}[label=\textup{(\alph*)},leftmargin=*]
\item (lifts) $|H^X(x',z)-E^S(x',z)|\le C(C^S_EJ_1\eta)^{1/2}e^{-\delta|x-z|/2}$ and $|H^Y(x'',z)-E^S(x',z)|\le C(375a\,J_1\eta)^{1/2}e^{-\delta|x-z|/2}$; the same for the photon lifts and for $E^S_j$ at level $j$ (unit scale of level $j$, $\eta\to\eta_j=5^{1-j}$).
\item (slices) for $2\le j<r$: $|S^X_j(x',y')-\tilde S^S_j(x',y'')|\le C(\eta'/\ell_j)^{1/2}\ell_j^{-1}e^{-\delta|x'-y'|/\ell_j}$.
\item (slices, the derivative at the lattice end) for $2\le j<r$: $|\partial_{y''}S^Y_j(x'',y'')-\partial_{y''}\tilde S^S_j(x',y'')|\le C(\eta'/\ell_j)^{1/2}\ell_j^{-2}e^{-\delta|x'-y'|/\ell_j}$.
\end{enumerate}
For $j\in\{0,1\}$ each term is bounded separately by \Cref{rot:hyb:lem:ESB} and the per-flow bounds of \Cref{rot:cor:unif}(a); there is no gain ($\ell_j\le5\eta'$).
\end{corollary}
\begin{proof}
(a) \Cref{rot:hyb:lem:ADMS} gives the sup bound $C^S_EJ_1\eta$; the second form is the same computation with $M^Y(q)$ at the cube offset $r_Y=x''-x$ in place of $M^X$: $|w_0u_Y-e^{\mathrm iq\cdot r_Y}u_Y|\le(0.106+c_\tau/5)\eta|q|$ and $\sum_{l\ne0}|w_l-e^{\mathrm i(q+l)\cdot r_Y}|\le124(W^3+1)$, giving $\le375a\,J_1\eta$. The decay is \Cref{rot:hyb:lem:ESB}(a) for $E^S$ and the separate bounds $|H^F(x^F,z)|\le Ce^{-\delta|x^F-z|}$ of \Cref{rot:lem:H}; take the geometric mean ($|x'-x|\le\sqrt{13}\eta'$). (b) In the unit scale of level $j$,
\[
S^X_j(x',y')-\tilde S^S_j(x',y'')=\sum_{z,z'}\big[(H^X_j-E^S_j)(x',z)\,C^X_j(z,z')\overline{H^X_j(y',z')}+E^S_j(x',z)(C^X_j-C^Y_j)(z,z')\overline{H^X_j(y',z')}+E^S_j(x',z)C^Y_j(z,z')\overline{(H^X_j(y',z')-H^Y_j(y'',z'))}\big]:
\]
(a), the bound $|C^X_j-C^Y_j|\le C\eta_je^{-\delta|z-z'|/2}$ of the proof of \Cref{rot:lem:slice}, \Cref{rot:hyb:lem:ESB}(a) and the decaying form of \Cref{rot:lem:H} at matched lattice points; sum the exponentials and rescale ($\eta_j=5\eta'/\ell_j$). This is \Cref{rot:lem:slice} with $E^S_j$ in place of $H^Y_j\circ\tau$ in the first argument. (c) The difference is $\sum_{z,z'}(H^Y_j(x'',z)-E^S_j(x',z))C^Y_j(z,z')\overline{\partial H^Y_j(y'',z')}$: the second form of (a), \cite[(3.6)]{K1}, and \cite[(3.7)]{K1} for the lattice gradient at $y''$; one more factor $\ell_j^{-1}$ on rescaling. (Alternatively, by \Cref{rot:hyb:lem:ESB}(c) the difference is $\cA(x'')-\cA(x')$ for $\cA\defeq\partial_{y''}\tilde S^S_j(\cdot,y'')$, so it is $\le c_\tau\eta'\sup|\nabla\cA|\le C(\eta'/\ell_j)\ell_j^{-2}e^{-\delta|x'-y'|/\ell_j}$ by \Cref{rot:hyb:lem:ESB}(b).)
\end{proof}
